\FloatBarrier

\section{Discussion and Conclusions}
\label{sec:discussion}

Reliable separation of unresolved and resolved sources is important for a
wide range of Rubin science applications.  Stellar-population studies require
samples with controlled galaxy contamination, while extragalactic analyses
require the complementary suppression of foreground stars.  These demands
become increasingly difficult to satisfy at faint magnitudes, where the
stellar and galaxy morphology distributions overlap and galaxies dominate the
source counts.  A continuous probabilistic description is therefore more
informative than a binary morphology flag because it retains the uncertainty
of each classification and allows independent measurements to be combined
before a catalog assignment is made
\citep{Henrion2011BayesianSG,Slater2020}.

XXXXXThe results in Section~\ref{sec:single-band-performance} show that most of the
improvement in the $r$-band ranking is obtained by modeling the full
magnitude-dependent morphology distributions.  For the 193,516 objects with a
finite $r$-band morphology measurement, the AUC increases from 0.871 for a
threshold sweep over the raw
$\Delta m_r=m_{{\rm PSF},r}-m_{{\rm CModel},r}$ measurement to 0.899 for the
morphology-only probabilistic score.  Including the source-count prior changes
the AUC to 0.901.  The comparatively small change in AUC does not make the
prior unimportant: AUC measures only the ordering of sources, whereas the
prior supplies the population odds required to interpret the morphology
likelihoods as posterior probabilities.  Its influence is greatest where the
morphology measurement is ambiguous and the two classes have very different
abundances \citep{Fawcett2006ROC}.

The fixed catalog assignment exposes the corresponding completeness--purity
trade-off.  At $16\leq r<24$, the probabilistic star-candidate sample has a
completeness of 0.878 and a purity of 0.975.  At $24\leq r<25$, its purity is
0.558, compared with 0.318 for the default Rubin
\texttt{extendedness} selection, while its completeness is 0.268, compared
with 0.582 for \texttt{extendedness}.  The probabilistic selection is thus
cleaner but less complete in this interval; it is not uniformly superior for
all scientific uses.  Analyses that prioritize rare-object discovery may
prefer higher completeness, whereas studies that are sensitive to galaxy
contamination may prefer the cleaner selection.  Reporting the continuous
$p_{\rm S}$ value preserves classification uncertainty in addition to the
adopted catalog assignment.  This class-specific reporting is particularly
important for the strongly imbalanced faint sample
\citep{SaitoRehmsmeier2015}.

At $25\leq r<26$, the adopted catalog assignment selects no star-like objects
in the evaluated sample.  The stellar completeness is consequently zero and
the purity of the empty star-candidate sample is undefined.  This result does
not imply that the field contains no stars.  Rather, the available morphology
information and the strongly galaxy-dominated source-count prior do not
produce posterior probabilities above the adopted assignment boundary.  The
continuous score still contains ranking information, but the present catalog
assignment is too conservative to define a useful faint stellar sample.  This
faint-end behavior is an important limitation rather than an improvement over
the Rubin baseline.

XXXXXThe multiband comparison leads to a similarly qualified conclusion.
Combining morphology measurements from additional bands improves the ranking
in the bright and intermediate intervals, but the gain is not monotonic with
the number of bands.  In the faintest interval, the $r$-only morphology score
slightly outperforms the six-band morphology score.  The color-augmented score
shown in Figure~\ref{fig:multiband-color-roc} improves the ranking relative to
the morphology-only curves within the evaluated sample, with its largest
displayed gain at the faint end.  Because the combined morphology-plus-color
quantity is used as a discriminative score rather than as a calibrated
likelihood ratio, it should be interpreted as a ranking statistic and not as
a posterior probability.  Likewise, the color distributions in
Figure~\ref{fig:color-distributions} describe the samples produced by the
classifier; because those colors contribute to the classification, that
figure is not an independent validation.  Previous work has similarly found
that multiband photometry can complement morphology, while also emphasizing
the dependence of performance on depth, labels, and survey conditions
\citep{Fadely2012MultiBandSG,Kim2015HybridSG,SevillaNoarbe2018DESY1}.

Several limitations determine how broadly the present measurements can be
generalized.  First, the validation is confined to the COSMOS field.  The
star-to-galaxy number ratio varies with Galactic coordinates, extinction,
depth, and catalog selection, so a prior inferred in COSMOS should not be
transferred unchanged across the Rubin footprint.  A spatially portable
version of the classifier could use a line-of-sight stellar-count prediction,
for example from TRILEGAL, together with an independently specified galaxy
count model \citep{Girardi2005TRILEGAL,DalTio2022TRILEGAL}.  TRILEGAL predicts
stellar counts rather than the galaxy-to-star ratio itself, and the stellar
and galaxy terms must be evaluated with consistent filters, extinction
conventions, magnitude bins, effective area, and catalog selections.

XXXXXSecond, 193,586 of the 269,606 selected Rubin sources are matched to the
COSMOS2020/FARMER reference catalog, a Rubin-side matching fraction of 0.718.
The reported metrics therefore characterize the matched validation sample,
not every selected Rubin source \citep{Weaver2022}.  The multiband comparisons
are further restricted to objects with finite values for every score included
in a given panel.  If matching failures or missing measurements depend on
class, magnitude, color, blending, or surface brightness, the evaluated sample
need not be representative of the full catalog.

Third, the external classifications are reference labels rather than
error-free physical truth.  High-resolution HST/ACS studies of COSMOS
demonstrate both the value and the magnitude dependence of morphology-based
reference samples \citep{Leauthaud2007COSMOSACS,Robin2007COSMOSStars}; we do
not assume that the COSMOS2020/FARMER labels used here are error free.
Independent checks with spectroscopic classifications and additional
high-resolution fields would help separate classifier errors from
reference-label errors and cosmic variance.

XXXXXFinally, the faint-end unresolved fraction is obtained by extrapolating a
smooth model from magnitude intervals in which the mixture components are
better separated.  The uncertainty of that extrapolation is not propagated
into the reported $p_{\rm S}$ values.  Future applications should propagate
uncertainties in the mixture parameters and population prior, and should test
sensitivity to depth, seeing, PSF modeling, blending, and crowding.  It is also
important to preserve the distinction between morphology and astrophysical
identity: $p_{\rm S}$ is formally the probability that a source belongs to the
unresolved morphological component.  Unresolved quasars and compact galaxies,
as well as resolved stellar systems, prevent this probability from being
universally identical to the probability that an object is physically a star.

XXXXXIn summary, modeling the full $\Delta m$ distribution extracts more ranking
information from the $r$-band morphology measurement than the binary Rubin
\texttt{extendedness} field and provides a continuous output that can be
combined with additional measurements.  In the matched COSMOS sample, the
probabilistic assignment produces a purer but less complete star-candidate
sample at $24\leq r<25$, while it becomes too conservative to select a useful
stellar sample at $25\leq r<26$.  Multiband morphology provides useful but
non-monotonic gains, and the color-augmented result is best regarded as a
ranking until it is probability-calibrated.  Extending the method across the
Rubin footprint will require field-dependent count priors, uncertainty
propagation, and validation in samples beyond the present COSMOS
complete-case data.